\par\Needspace{6\baselineskip}\medskip
\SourceLeaf{099}\hypertarget{native-099-article075}{}\pdfbookmark[2]{75.}{article-075}
\begin{center}75.\end{center}
Another solution to the preceding problem rests on the observation in Article 64, III that the heliocentric position of the Earth, the geocentric position of the celestial body, and its heliocentric position all lie on one and the same great circle of the sphere. Let these positions in Fig.~3 be $T$, $G$, and $H$, respectively. Let $\AscendingNode$ be the position of the ascending node, $\AscendingNode T$ and $\AscendingNode H$ parts of the ecliptic and the orbit, and $GP$ the perpendicular drawn from $G$ to the ecliptic, which will therefore be $=b$. From this and the arc $PT=L-l$ we determine the angle $T$ and the arc $TG$. In the spherical triangle $\AscendingNode HT$ we are then given the angle $\AscendingNode=i$, the angle $T$, and the side $\AscendingNode T=L-\AscendingNode$, from which we find the two remaining sides $\AscendingNode H=u$ and $TH$. Finally we have $HG=TG-TH$ and $r=\frac{R\sin TG}{\sin HG}$, $\Delta=\frac{R\sin TH}{\sin HG}\mathpunct{\cdot}$
